%% kapa: exam=1 round=33 year=2022 period=2 form=A subject=accounting subject_ko=회계학 seq=11 number=51
%% source: Q-Net 감정평가사 기출문제 (cst00309) artlSeq=5209997 file=raw/1008/5209997/2022년도 제33회 감정평가사 1차시험 2교시(감정평가 관계법규, 회계학).pdf
%% conversion: kordoc 4.24.51 -> tools/gen1.py (auto); review=proofread-sonnet; flags=none
\begin{question}{51}
다음은 (주)감평의 20×1년도 재무제표의 일부 자료이다.
\begin{framedtext}
(1) 재무상태표의 일부 자료 \par
{\small\centering\begin{tabular}{|c|r|r|}\hline
계정과목 & 기초잔액 & 기말잔액 \\ \hline
매출채권(순액) & ￦140 & ￦210 \\
선급영업비용 & 25 & 10 \\
미지급영업비용 & 30 & 50 \\ \hline
\end{tabular}\par}
(2) 포괄손익계산서의 일부 자료 \par
{\small\centering\begin{tabular}{|c|r|}\hline
매 출 액 & ￦410 \\
영 업 비 용 & 150 \\ \hline
\end{tabular}\par}
\end{framedtext}
위 자료에 기초한 20×1년도 (주)감평의 (A)고객으로부터 유입된 현금흐름과 (B)영업비용으로 유출된 현금흐름은?
\begin{choices}
  \item (A) ￦335 \quad (B) ￦155
  \item (A) ￦340 \quad (B) ￦115
  \item (A) ￦340 \quad (B) ￦145
  \item (A) ￦350 \quad (B) ￦115
  \item (A) ￦350 \quad (B) ￦155
\end{choices}
\end{question}
\answer{2}
